
\label{subsection:model}
\newcommand{\comDAG}{\pi_{C[\mathcal{E}]}} 
\paragraph{The Computation Model.} Informally, a quantum circuit is a placement of a collection of quantum gates, taken from a finite set, at different space-time locations. Each such location records when the gate is applied and on which qubits it acts. Formally:

\begin{definition}[Quantum Circuit]
  \label{def:quantcirc}
  A \emph{space-time location} is a pair $l=(t,S)$, where $t$ is a time coordinate and $S=(s_{1},\ldots,s_{\Delta})$ is an ordered tuple of qubit indices. For a finite collection of quantum gates
  \begin{equation*}
    \begin{split}
      \mathcal{G} \subset \Big\{ g \colon L(\mathcal{H}_{2}^{\otimes \Delta}) \to L(\mathcal{H}_{2}^{\otimes \Delta}) \Big\},
    \end{split}
  \end{equation*}
  and integers $w,d\in\mathbb{N}$, a quantum circuit $C$ of depth $d$ and width $w$ is a collection of tuples
  \begin{equation*}
    \begin{split}
      C=\{u_i=(g_i,l_i)\}_i,
    \end{split}
  \end{equation*}
  where each $g_i\in\mathcal{G}$ and each location $l_i=(l_{i,1},l_{i,2})$ belongs to $[d]\times [w]^{\Delta}$. We call the tuples $u_i$ \textit{the gates of $C$}. They must satisfy the following compatibility condition: if two gates have the same time coordinate, then their supports are disjoint. Namely, for $u_i=(g_i,l_i),u_j=(g_j,l_j)\in C$, the equality $l_{i,1}=l_{j,1}$ implies that the sets of qubit indices appearing in $l_{i,2}$ and $l_{j,2}$ are disjoint. We denote by
  \begin{equation*}
    \begin{split}
      \operatorname{Loc}(C) \colonequals \{l_i : u_i=(g_i,l_i)\in C\}
    \end{split}
  \end{equation*}
  the set of locations occupied by the gates of $C$. We define the \textit{$t$-layer} of $C$ to be its restriction to gates at time $t$:
  \begin{equation*}
    \begin{split}
      C|_{t} \colonequals \{u_i=(g_i,l_i)\colon u_i\in C \text{ and } l_{i,1}=t\}.
    \end{split}
  \end{equation*}
\end{definition}
Later, we will restrict $\mathcal{G}$ to contain $\mathbf{CX}, \mathbf{H}, \mathbf{S}, \mathbf{T}$, and measurements. To define the action of a quantum circuit $C$ on a given mixed state, we first associate $C$ with a directed acyclic graph.
\begin{definition}[Associated Computation DAG]
  \label{def:steps}
  Let $C=\{u_i\}_i$ be a quantum circuit with depth $d$ and width $w$. We define the associated computation DAG of $C$, denoted by $\pi_C$, to be the directed acyclic graph $\pi_C=(\mathcal{U},E)$ as follows:
  \begin{enumerate}
    \item The vertex set is the gate set of $C$, namely $\mathcal{U}=\{u_i\}_i$.
    \item For $u_i=(g_i,l_i)$ and $u_j=(g_j,l_j)$, there is a directed edge $(u_i,u_j)\in E$ if and only if $j$ is one of the minimizers of
      \begin{equation*}
        \begin{split}
          \big\{ l_{k,1} : (g_k,l_k)\in \mathcal{U},\ l_{k,1}>l_{i,1},\ l_{k,2}\cap l_{i,2}\neq\emptyset \big\}.
        \end{split}
      \end{equation*}
      Equivalently, $u_j$ is one of the earliest later gates whose support intersects the support of $u_i$. There may be more than one such minimizer, so the out-degree may exceed one.
  \end{enumerate}

  Suppose that $J=(j_1,j_2,\ldots,j_k)$ is an ordering of the gate indices that agrees with a topological order induced by $\pi_C$. Then we define the associated partial circuits $C_J=\{C_{J,i}\}_{i=1}^{k}$ by
  \begin{enumerate}
    \item $C_{J,1}=\{u_{j_1}\}$.
    \item For $1\le i<k$, set $ C_{J,i+1}=C_{J,i}\cup\{u_{j_{i+1}}\}$.
  \end{enumerate}
      We call $C_{J,i}$ the partial computation circuit up to step $i$.
\end{definition}

\begin{definition}
  Let $\rho$ be a density matrix over $w$ qubits, and let $C = \{u_{i}\}_{i=1}$ be a quantum circuit with depth $d$ and width $w$. For any density matrix $\sigma$, quantum gate $g \in \mathcal{G}$, and location $l$, the application of the gate-location tuple $(g,l)$ on $\sigma$, denoted by $(g,l) \circ \sigma$, is the application of $g$ on $\sigma$ where wires are ordered such that the first qubit in the input to $g$ is $l_{2,1}$, the second is $l_{2,2}$, and so on. The result of the application of $C$ on $\rho$, denoted by $C \circ \rho$, is given by iterative applications of the gates $\{u_{i}\}_{i=1}$ on $\rho$, according to a topological order induced by $\pi_{C}$.

We say that $\rho_{1}, \rho_{2}, \ldots, \rho_{k}$ is a computation prefix if $\rho_{1} = \rho$ and $\rho_{i+1} = u^\prime_{i} \circ \rho_{i}$, where $\{u^{\prime}_{i}\}_{i=1}$ is a prefix of a topological ordering of $\{u_{i}\}_{i=1}$. If $\{u^{\prime}_{i}\}_{i=1}$ contains all the gates in $\{u_{i}\}_{i=1}$, then we call $\rho_{1}, \rho_{2}, \ldots, \rho_{k}$ a running.
\end{definition}
For a quantum circuit $C$, we would like to discuss its action in the presence of different outcomes of faults. The following definition encodes the running in the presence of a specific set of faults $\mathcal{E}$.
%We do this by defining a new circuit, stacking a layer of $\mathcal{E}$ after every $m$ layers of $C$.
\begin{definition}[$C$ subjected to $\mathcal{E}$ every $m$ computational layers]
  Let $C=\{u_i\}_i$ be a quantum circuit with depth $d$ and width $w$, let $m\in\mathbb{N}$, and set
  \begin{equation*}
    \begin{split}
      d_m \colonequals \left\lceil \frac{d}{m} \right\rceil.
    \end{split}
  \end{equation*}
  Let $\mathcal{L}\subset [d_m]\times [w]$ be a set of locations, and let $\mathcal{E}$ be a quantum circuit, such $\text{Loc}(\mathcal{E}) = \mathcal{L}$. We name $\mathcal{E}$ \textit{subsets of faults}, and we name the layers of $\mathcal{E}$ \textit{fault layers}.
  The circuit \textit{$C$ subjected to $\mathcal{E}$ every $m$ computational layers}, denoted by $C[\mathcal{E}]_{m}$, is obtained by setting the first layer to be the first layer of $\mathcal{E}$, then placing in order the layers of $C$ and inserting one fault layer after every $m$ consecutive $C$'s layers. Thus the second fault layer is inserted after the layers $1,\ldots,m$ of $C$, the third is inserted after the layers $m+1,\ldots,2m$ of $C$, and so on. Concretely, each gate $u_i=(g_i,(t_i,S_i))$ of $C$ is moved to the time coordinate
  \begin{equation*}
    \begin{split}
      t_i + \left\lfloor \frac{t_i-1}{m} \right\rfloor,
    \end{split}
  \end{equation*}
  while each fault gate $(g,(t,S))\in\mathcal{E}$ is placed at time coordinate
  \begin{equation*}
    \begin{split}
      t(m+1).
    \end{split}
  \end{equation*}
  When $m=1$, the gates occupy the even time coordinates and the fault layers occupy the odd time coordinates. When the parameter is omitted, we mean $m=1$. We denote the associated computation DAG of $C[\mathcal{E}]_{m}$ by $\comDAG$.
\end{definition}





\begin{remark}
  Formally, a fault circuit is just a particular quantum circuit whose gates are interpreted as faults rather than as computational operations. Thus the distinction between a quantum circuit and a fault circuit is semantic rather than structural: both are circuits in the sense of \Cref{def:quantcirc}, but they play different roles in the analysis.
\end{remark}

\begin{example}
  When $m=3$, the 'computational' layers are sent to the times
  \begin{equation*}
    \begin{split}
      2,3,4,6,7,8,10,11,12,\ldots,
    \end{split}
  \end{equation*}
  while the fault layers are sent to the times
  \begin{equation*}
    \begin{split}
      1,5,9,13,\ldots.
    \end{split}
  \end{equation*}
\end{example}


\begin{definition}[Noise channel and local stochastic noise channel]
  \label{def:localstoc}
  A \textit{noise channel} $\mathcal{N}_{\mathcal{E}}$ is a distribution over sets of fault locations. Let $p\in(0,1)$, we say that the noise channel $\mathcal{N}_{\mathcal{E}}$ is \textit{local stochastic with rate $p$} if for every finite set of locations $S\subset \mathbb{N}\times \mathbb{N}$,
  \begin{equation*}
    \begin{split}
      \prbm{S\subset \operatorname{Loc}(\mathcal{E})}{\mathcal{E}\sim \mathcal{N}_{\mathcal{E}}}
      \le p^{|S|}.
    \end{split}
  \end{equation*}
\end{definition}
So, $N_{\mathcal{E}}$ defines a probability space. The set of faults $\mathcal{E}$ is an atomic outcome, and for a given circuit $C$, the subjected circuit $C[\mathcal{E}]$ is a random variable. The following claim states that for local stochastic noise, the subjection of every 1-layer is equivalent to the subjection of every $m$-layer, up to a rescaling of the noise rate. Here, by equivalence, we mean that the actions of the subjected circuits are the same\footnote{For a given input (mixed) state $\rho$, they yield the same output}. 

\inb{\textbf{Rewrite.}} We give the proof in the appendix for several reasons. First, the argument is not the core of this work. Second, versions of it appear in almost any fault tolerance work. Additionally, it relies on 'propagation of entire layers,' a term that conflicts with the 'propagation to a step' (\Cref{aa}) needed for the intermediate theorem.
\begin{claim}
  \label{claim:block_noise_renorm}
  Let $\mathcal{N}_{\mathcal{E}}$ be a local stochastic noise channel with rate $p$, and let $C$ be a quantum circuit whose gates have arity at most $\Delta$. Fix an even integer $c$. Then there exists a local stochastic noise channel $\mathcal{N}_{\mathcal{E}_{q}}$ with rate
  \begin{equation*}
    \begin{split}
      q \colonequals \left(2^{H(\Delta^{-2c})\Delta^{2c}}p\right)^{\Delta^{-c}},
    \end{split}
  \end{equation*}
  such that for any given state $\rho$ it holds that:
  \begin{equation*}
    \begin{split}
      C[\mathcal{E}]_{1} \circ \rho =  C[\mathcal{E}_{q}]_{c/2} \circ \rho
    \end{split}
  \end{equation*}
  In particular, up to a constant depending only on $\Delta$ and $c$, the rate $q$ behaves like $p^{\Delta^{-c}}$.
\end{claim}
Equipped with this claim, we are allowed to assume having constant depth cycles that don't observe faults. 


\begin{definition}
  Let $C$ be a quantum circuit that contains $l$ measurements. For possible outcomes $\alpha \in \mathbb{F}^{l}_{2}$, let $C[,\alpha]$ be the circuit obtained by conditioning that $\alpha_{1},..,\alpha_{l}$ were measured. We extend the notation for the subjected circuit by denoting $C[\mathcal{E},\alpha]$ as the subjected circuit where the measurements are assumed to outcome $\alpha_{1},..,\alpha_{l}$.
\end{definition}
In our case, the conditioned circuit $C[\mathcal{E}, \alpha]$ is a Clifford circuit; in particular, it sends Paulis to Paulis. That brings us to define, for a given set of faults $\mathcal{E}$ and outcomes $\alpha$, a Pauli-path that describes the evolution of the faults through the circuit.
\inb{\textbf{ For that being immediate we have to measure the stabilizer after every noise layer. Also,  }} 

\begin{definition}[Propagation]
  \label{def:propagation}
Let $C$ be a quantum circuit at depth $\mu$, and let $\alpha \in \mathbb{F}_{2}^{(\cdot)}$ be a possible outcome vector for $C$. Suppose that $C$ is a Clifford+measurements circuit, so $C[,\alpha]$ is a Clifford circuit. We define the propagation of Pauli $P \in \mathcal{P}$ through the conditioned circuit $C[,\alpha]$ to be the Paulis $P_{1}, \ldots, P_{\mu+1}$ such that $P_{1} = P$ and for any $i \in [2, \mu+1]$:
  \begin{equation*}
    \begin{split}
      P_{i}C[,\alpha]|_{\rightarrow i} =  C[,\alpha]|_{ i} P_{i-1} C[,\alpha]|_{\rightarrow i-1} = C[,\alpha]|_{\rightarrow i}P 
    \end{split}
  \end{equation*}
We extend the notation to step-propagation by considering the propagation of the circuit obtained from $C[\mathcal{E},D]$ by flattening the circuit such that each gate is applied at the time that equals its step. 
\end{definition}


